\documentclass[a4paper]{article}

% Español (México) con XeLaTeX: fontspec para fuentes Unicode, polyglossia
% para la división silábica y los nombres del documento en la variante mexicana.
\usepackage{fontspec}
\usepackage{polyglossia}
\setdefaultlanguage[variant=mexican]{spanish}
% Los nombres chinos van como «pinyin (original)»: xeCJK da a los caracteres
% Han su propia fuente; sin ella XeLaTeX los omite con un aviso.
\usepackage{xeCJK}
\setCJKmainfont{FandolSong-Regular.otf}
% polyglossia no traduce los nombres de listings.
\AtBeginDocument{\providecommand{\lstlistingname}{}\renewcommand{\lstlistingname}{Listado}%
  \providecommand{\lstlistlistingname}{}\renewcommand{\lstlistlistingname}{Índice de listados}}
\usepackage{amsmath, amssymb}
\usepackage{graphicx}
\usepackage[margin=2.5cm]{geometry}
\usepackage[most]{tcolorbox}
\usepackage{etoolbox}
\usepackage{listings}
\usepackage{hyperref}
\usepackage{booktabs}
\usepackage{subcaption}
\usepackage{float}

\newtcolorbox{knowledgebox}[1]{
    enhanced, colback=blue!5!white, colframe=blue!75!black, colbacktitle=blue!75!black,
    coltitle=white, fonttitle=\bfseries, title={#1}, attach boxed title to top left={yshift=-2mm, xshift=2mm},
    boxrule=1pt, sharp corners
}

\newtcolorbox{importantbox}[1]{
    enhanced, colback=yellow!10!white, colframe=yellow!80!black, colbacktitle=yellow!80!black,
    coltitle=black, fonttitle=\bfseries, title={#1}, sharp corners
}

\newtcolorbox{warningbox}[1]{
    enhanced, colback=red!5!white, colframe=red!75!black, colbacktitle=red!75!black,
    coltitle=white, fonttitle=\bfseries, title={#1}, sharp corners
}

\lstset{
    language=Python,
    basicstyle=\ttfamily\small,
    keywordstyle=\color{blue},
    stringstyle=\color{red!60!black},
    commentstyle=\color{green!60!black},
    breaklines=true,
    frame=single,
    numbers=left,
    numberstyle=\tiny\color{gray},
    captionpos=b,
    extendedchars=true
}

\newcommand{\notetitle}{MoE, primeros pasos: cálculo del número de parámetros, Qwen3-30B-A3B, GQA y Sparse MoE}
\newcommand{\noteauthors}{Elaboradas a partir de materiales públicos del curso}
\newcommand{\notedate}{2025}
\newcommand{\videochannel}{Wudaokou Naishi (五道口纳什)}
\newcommand{\videopublishdate}{2025}
\newcommand{\videoduration}{aproximadamente 15 minutos}
\newcommand{\videourl}{}
\newcommand{\videocoverpath}{cover.jpg}
\newcommand{\repourl}{https://github.com/hqhq1025/ai-course-notes}


% Footer with repo info
\usepackage{fancyhdr}
\pagestyle{fancy}
\fancyhf{}
\fancyfoot[C]{\thepage}
\fancyfoot[R]{\footnotesize\color{gray}\href{\repourl}{AI Course Notes}}
\renewcommand{\headrulewidth}{0pt}
\renewcommand{\footrulewidth}{0pt}

\begin{document}

\begin{titlepage}
\centering
{\Large Notas del curso\par}
\vspace{1.2cm}
{\huge\bfseries \notetitle\par}
\vspace{0.8cm}
{\large \noteauthors\par}
\vspace{0.3cm}
{\large \notedate\par}
\vspace{1.2cm}

\ifdefempty{\videocoverpath}{
{\small Al generar las notas, indique la ruta local de la portada del video.\par}
}{
\includegraphics[width=0.82\textwidth,height=0.45\textheight,keepaspectratio]{\videocoverpath}\par
}

\vfill
\begin{tcolorbox}[width=0.9\textwidth, colback=black!2!white, colframe=black!60, sharp corners]
\textbf{Autor o canal del video}: \videochannel\par
\textbf{Fecha de publicación}: \videopublishdate\par
\textbf{Duración del video}: \videoduration\par
\ifdefempty{\videourl}{}{
\textbf{Enlace del video}: \href{\videourl}{\nolinkurl{\videourl}}\par
}
\par
\textbf{Página del proyecto}: \href{\repourl}{\nolinkurl{\repourl}}\par
\end{tcolorbox}
\end{titlepage}

\tableofcontents
\newpage

%% --- Inicio del contenido --- %%

\section{Introducción: ¿por qué prestar atención a MoE?}

Esta es la primera clase de la serie «Arquitectura de modelos de lenguaje grandes». En los últimos años, la evolución de las arquitecturas de los modelos grandes ha sido vertiginosa, y entre ellas la arquitectura \textbf{Mixture of Experts (MoE)} ha atraído mucha atención por su enorme ventaja en eficiencia de inferencia.

Tomemos como ejemplo DeepSeek R1: es un modelo MoE con alrededor de 670,000 millones de parámetros totales, pero en la inferencia solo activa alrededor de 72,000 millones de parámetros. En comparación con un modelo Dense completo de 72,000 millones de parámetros, la velocidad de inferencia de ambos es básicamente la misma; esto significa que un modelo MoE, al mantener una capacidad de modelo enorme (gran capacidad), logra una eficiencia de inferencia muy alta (pocos parámetros activados).

\begin{importantbox}{El valor central de MoE}
\begin{itemize}
  \item \textbf{Un número total de parámetros grande} $\Rightarrow$ gran capacidad del modelo, gran capacidad de desempeño
  \item \textbf{Un número de parámetros activados pequeño} $\Rightarrow$ inferencia rápida, alta eficiencia
  \item En la época de Qwen2.5 todavía se usaban modelos totalmente Dense; para Qwen3 ya se adoptan ampliamente variantes MoE
\end{itemize}
\end{importantbox}

En esta clase se toman \textbf{Qwen3-32B} (Dense) y \textbf{Qwen3-30B-A3B} (MoE) como objeto de estudio, y a través de comparar y calcular su número de parámetros se busca entender la arquitectura MoE.
\section{Dense vs. MoE: comparación de configuraciones}

\subsection{Panorama de configuraciones de modelos}

\begin{table}[H]
\centering
\caption{Comparación de configuraciones entre Qwen3-32B (Dense) y Qwen3-30B-A3B (MoE)}
\begin{tabular}{lcc}
\toprule
\textbf{Elemento de configuración} & \textbf{Qwen3-32B (Dense)} & \textbf{Qwen3-30B-A3B (MoE)} \\
\midrule
Número de capas (Layers) & 64 & 48 \\
Dimensión del modelo $d_{\text{model}}$ & 5120 & 2048 \\
Número de cabezas de atención & 64 & 32 \\
Head Dimension & 128 & 128 \\
Número de grupos KV (GQA) & 8 & 4 \\
Dimensión intermedia de la FFN & 25600 & 768 (por Expert) \\
Número total de Experts & — & 128 \\
Número de Experts activados & — & 8 \\
Tie Word Embeddings & False & False \\
\bottomrule
\end{tabular}
\end{table}

\begin{knowledgebox}{¿Por qué los modelos MoE son «más angostos y más superficiales»?}
Del Dense de 64 capas al MoE de 48 capas, $d_{\text{model}}$ pasa de 5120 a 2048 y el número de cabezas de atención de 64 a 32. El modelo MoE compensa estas reducciones de dimensión mediante una gran cantidad de experts (128), y el número total de parámetros sigue siendo grande ($\sim$30B), pero en cada inferencia solo se activan 8 experts, con lo cual el número de parámetros activados es de apenas unos 3.3B.
\end{knowledgebox}

\subsection{Línea de evolución de la arquitectura}

La mayoría de las arquitecturas de LLM dominantes actuales se derivan de la ruta técnica Transformer $\to$ LLaMA 2:

\begin{itemize}
  \item \textbf{RMSNorm}: sustituye a LayerNorm, elimina el centrado de la media y solo aplica el reescalamiento
  \item \textbf{SwiGLU}: sustituye a ReLU como función de activación de la FFN
  \item \textbf{RoPE}: codificación posicional rotativa
  \item \textbf{GQA}: atención de consultas agrupadas, que comprime la KV Cache
\end{itemize}

El cambio central de la arquitectura MoE consiste en \textbf{sustituir el módulo de FFN de cada capa por un módulo Sparse MoE}, es decir, pasar de un MLP compartido de dos capas a 128 FFN pequeñas e independientes (Experts), de las cuales en cada paso solo se enruta hacia las top-8.

\section{GQA (Grouped Query Attention) explicado en detalle}

\subsection{Motivación y mecanismo de GQA}

\begin{importantbox}{Idea central de GQA}
En la MHA estándar, cada cabeza de Query corresponde a una cabeza de KV independiente. GQA hace que varias cabezas de Query \textbf{compartan} un mismo grupo de cabezas de KV, con lo cual se reduce considerablemente el costo de almacenamiento de la KV Cache.
\end{importantbox}

Tomando como ejemplo Qwen3-30B-A3B:
\begin{itemize}
  \item 32 cabezas de Query, 4 grupos de KV
  \item Cada $32 / 4 = 8$ cabezas de Query comparten un grupo de $K, V$
  \item Cabezas de Query 0--7 $\to$ grupo de KV 0; cabezas de Query 8--15 $\to$ grupo de KV 1; y así sucesivamente
\end{itemize}

\subsection{$d_{\text{model}} \neq n_{\text{heads}} \times d_{\text{head}}$}

\begin{warningbox}{Las dimensiones no siempre coinciden}
En LLaMA 2 o GPT-2, $d_{\text{model}} = n_{\text{heads}} \times d_{\text{head}}$. Pero en el modelo Dense de Qwen3, $d_{\text{model}} = 5120$, mientras que $n_{\text{heads}} \times d_{\text{head}} = 64 \times 128 = 8192 \neq 5120$. Esto significa que $W_Q$ ya no es una matriz cuadrada, sino una matriz no cuadrada de $5120 \times 8192$.
\end{warningbox}

\subsection{Resumen de la sección}

GQA comprime la KV Cache haciendo que varias cabezas de Query compartan grupos de KV, y es un estándar en los LLM predominantes actuales. Qwen3 va más allá y rompe la restricción de alineación entre $d_{\text{model}}$ y la dimensión de las cabezas, usando matrices de proyección no cuadradas.

\section{Arquitectura Sparse MoE en detalle}

\subsection{Estructura de la capa MoE}

Cada capa MoE contiene los siguientes componentes:

\begin{enumerate}
  \item \textbf{Router (enrutador)}: mapea el embedding de $d_{\text{model}}$ dimensiones de un token a un vector de logits de 128 dimensiones
  \item \textbf{Selección Top-K}: de los 128 Expert se eligen los top-8
  \item \textbf{Softmax Gating}: se aplica softmax a los logits de los top-8, obteniendo los coeficientes de ponderación
  \item \textbf{Expert FFN}: cada Expert es un MLP de dos capas ($W_{\text{up}}, W_{\text{gate}}, W_{\text{down}}$)
  \item \textbf{Suma ponderada}: las salidas de los 8 Expert se suman de forma lineal, ponderadas por los coeficientes de gating
\end{enumerate}

\subsection{Mecanismo del Router}

$$
\text{Router}: \mathbb{R}^{d_{\text{model}}} \to \mathbb{R}^{N_{\text{experts}}}
$$

Es decir, $W_{\text{router}} \in \mathbb{R}^{2048 \times 128}$, que mapea el embedding de cada token a los puntajes correspondientes a los 128 Expert.

Después de elegir los top-8, se aplica softmax a esos 8 puntajes (controlado por el parámetro de configuración \texttt{norm\_topk\_prob}, cuyo valor predeterminado es True), obteniendo los coeficientes de gating normalizados $g_1, g_2, \ldots, g_8$.

\subsection{Estructura de la FFN del Expert}

Cada Expert es un MLP de dos capas con SwiGLU:

$$
\text{Expert}(x) = W_{\text{down}} \cdot \big[ \text{SiLU}(W_{\text{up}} \cdot x) \odot (W_{\text{gate}} \cdot x) \big]
$$

donde:
\begin{itemize}
  \item $W_{\text{up}} \in \mathbb{R}^{2048 \times 768}$: mapea $d_{\text{model}}$ a la dimensión intermedia
  \item $W_{\text{gate}} \in \mathbb{R}^{2048 \times 768}$: la rama de gating, multiplicada elemento a elemento con la salida de $W_{\text{up}}$
  \item $W_{\text{down}} \in \mathbb{R}^{768 \times 2048}$: mapea la dimensión intermedia de vuelta a $d_{\text{model}}$
\end{itemize}

\begin{knowledgebox}{Dense FFN vs. MoE Expert FFN}
En el modelo Dense, la dimensión intermedia de la FFN es $25600 = 5 \times d_{\text{model}}$, mientras que en MoE la dimensión intermedia de cada Expert es apenas $768$. Aunque un solo Expert es mucho más pequeño que la FFN Dense, la cantidad total de parámetros de los 128 Expert es $128 \times 3 \times 2048 \times 768 \approx 604\text{M}$ (solo la parte de FFN); sumando las demás capas, el total de parámetros es de aproximadamente 30B.
\end{knowledgebox}

\subsection{Salida final de MoE}

$$
\text{MoE}(x) = \sum_{i=1}^{8} g_i \cdot \text{Expert}_i(x)
$$

Cada Expert produce un vector de 2048 dimensiones; estos se suman de forma ponderada según los coeficientes de gating para obtener la salida final de la capa MoE.

\subsection{Resumen de la sección}

La idea central de Sparse MoE es «enrutamiento + activación dispersa»: el Router elige, para cada token, los top-K Expert más relevantes, y el resto de los Expert no participan en el cómputo. Esto permite que el modelo tenga una cantidad total de parámetros muy grande mientras el cómputo de cada inferencia es reducido.
\section{Cálculo del número de parámetros}

\subsection{Token Embedding}

Qwen3-30B-A3B tiene \texttt{tie\_word\_embeddings = False}, es decir, los embedding de entrada y de salida \textbf{no se comparten}, por lo que se necesitan dos:

$$
P_{\text{embed}} = 2 \times V \times d_{\text{model}}
$$

donde $V$ es el tamaño del vocabulario y $d_{\text{model}} = 2048$.

\subsection{Número de parámetros por capa}

Cada capa contiene tres partes:

\textbf{1. Parte de atención GQA}:
$$
P_{\text{attn}} = d_{\text{model}} \times (n_Q \cdot d_h) + 2 \times d_{\text{model}} \times (n_{KV} \cdot d_h) + (n_Q \cdot d_h) \times d_{\text{model}}
$$

donde $n_Q = 32, n_{KV} = 4, d_h = 128$:
$$
P_{\text{attn}} = 2048 \times 4096 + 2 \times 2048 \times 512 + 4096 \times 2048
$$

\textbf{2. Router}:
$$
P_{\text{router}} = d_{\text{model}} \times N_{\text{experts}} = 2048 \times 128
$$

\textbf{3. Expert de MoE} (los 128 Expert en total):
$$
P_{\text{MoE}} = N_{\text{experts}} \times 3 \times d_{\text{model}} \times d_{\text{ffn}} = 128 \times 3 \times 2048 \times 768
$$

\subsection{Número total de parámetros y número de parámetros activados}

\begin{importantbox}{El origen de 30B frente a 3.3B}
\begin{itemize}
  \item \textbf{Número total de parámetros} ($\sim$30.5B): Token Embedding + 48 capas $\times$ (Attention + Router + \textbf{128} $\times$ Expert FFN)
  \item \textbf{Número de parámetros activados} ($\sim$3.3B): Token Embedding + 48 capas $\times$ (Attention + Router + \textbf{8} $\times$ Expert FFN)
  \item La única diferencia entre ambos: el número de Expert pasa de 128 $\to$ 8
\end{itemize}
\end{importantbox}

Esto es precisamente lo que significa el nombre del modelo «30B-A3B»: el número total de parámetros es de aproximadamente 30B, y el número de parámetros activados es de aproximadamente 3B.

\subsection{Resumen de la sección}

Mediante el cálculo término por término, vemos con claridad cómo se distribuye el número de parámetros del modelo MoE: una gran cantidad de parámetros se almacena en los 128 Expert, pero durante la inferencia solo se activan 8. Esto logra el objetivo de «gran capacidad, bajo cómputo».

\section{Síntesis y ampliación}

En esta sesión se partió de las dos variantes de Qwen3, Dense y MoE, y mediante la comparación de configuraciones y el cálculo del número de parámetros paso por paso se construyó una comprensión intuitiva de la arquitectura Sparse MoE:

\begin{enumerate}
  \item MoE sustituye la FFN por múltiples Expert independientes más un Router, lo que logra una activación dispersa
  \item GQA comprime la KV Cache compartiendo cabezas KV, y es hoy el estándar de facto en los LLM
  \item Qwen3 rompe la restricción tradicional $d_{\text{model}} = n_{\text{heads}} \times d_{\text{head}}$
  \item De los 30B de parámetros totales, en inferencia solo se activan alrededor de 3.3B, con una eficiencia de inferencia cercana a la de un modelo Dense del mismo tamaño
\end{enumerate}

\subsection{Lecturas adicionales}

\begin{itemize}
  \item El blog de Sebastien Raschka: serie de explicaciones detalladas sobre arquitecturas de LLM
  \item El informe técnico oficial de Qwen3
  \item El diseño de MoE en los informes técnicos de DeepSeek V2/V3
  \item Las explicaciones sobre LLaMA 2 y GQA en la serie «Post NOS Chat GPT»
\end{itemize}

%% --- Fin del contenido --- %%

\end{document}
